%%%%%%%%%%  scan idx 216  |  folio 203  %%%%%%%%%%
(Si $X$ est
paracompact et $\cO_X$ mou (cf. II 2.3.1) on peut compléter la suite exacte (1.2.4.2) en ajoutant un zéro à droite). L'homomorphisme $\sigma_n$ est appelé \emph{symbole}~; il associe à un opérateur différentiel de $E$ dans $F$ d'ordre $\le n$ une «~application polynomiale homogène de degré $n$ de $\Omega^1_{X/S}$ dans $\underline{\Homop}_{\cO_X}(E,F)$ (resp. $\underline{\Homop}_{\cO_X\otimes\mathbb C}(E,F)$)~», que l'on peut regarder encore, si l'on veut, comme un homomorphisme $p^*(E)\longrightarrow p^*(F)$ (où $p$ est la projection sur $X$ du fibré cotangent relatif (1.2.1)), satisfaisant une certaine condition d'homogénéité.

Soient $E,F,G$ des $\cO_X$-Modules localement libres de type fini, et $u\in\Diff^m(E,F)$, $v\in\Diff^n(F,G)$, d'où $vu\in\Diff^{m+n}(E,G)$. On vérifie facilement la relation
\begin{equation}\tag{1.2.4.3}
  \sigma_{m+n}(vu)=\sigma_n(v)\sigma_m(u).
\end{equation}
(On a la même relation pour les opérateurs différentiels à coefficients complexes). En particulier, si le composé de deux opérateurs différentiels est nul, le composé de leurs symboles l'est aussi.

\textbf{1.2.5.} Soit $E$ un complexe, à degrés bornés, d'opérateurs différentiels à coefficients complexes sur $X$ relativement à $S$ (1.2.3), $d^i:E^i\longrightarrow E^{i+1}$ étant d'ordre $\le n_i$. D'après (1.2.4.3), on a
\[
  \sigma_{n_{i+1}}(d^{i+1})\,\sigma_{n_i}(d^i)=0.
\]
On dit que $E$ est un \emph{complexe elliptique} d'ordre $\underline{n}=(n_i)$ si le complexe $(p^*(E^i),\sigma_{n_i}(d^i))$ (où $p$ est la projection sur $X$ du fibré cotangent) est acyclique en dehors de la section nulle.

Quand $S$ est réduit à un point, la notion précédente de complexe elliptique sur $X$ coïncide avec la notion usuelle, étudiée en détail dans [4] ou [16].

\textbf{1.2.6.} Soient $s\in S$, $i_s$ l'inclusion de la fibre $X_s$ dans $X$ (considérée comme monomorphisme d'espaces annelés) cf.~(1.1)). Pour $n\in\N$, on a un isomorphisme canonique
